\documentclass[11pt,letterpaper]{article}
\usepackage{amsmath,amssymb,fullpage,graphicx,pgf,float,hyperref}
\providecommand{\mathdefault}[1]{#1}
\title{42-Foot Fin-Keel Sailboat: Vessel RAO Simulation Charts}
\author{Mikhail Grushinskiy}
\date{Geometrically scaled RAO family}
\newcommand{\VesselChart}[2]{
  \begin{figure}[H]\centering
  \resizebox{\textwidth}{!}{\input{vessel_rao_42ft_#1.pgf}}
  \caption{#2}
  \end{figure}\clearpage
}
\begin{document}
\maketitle
These charts show vessel motion, rather than surface-particle motion, for an
estimated 42-foot fin-keel sailboat at zero forward speed with fixed heading
and an IMU at the center of gravity. Waterline length is 10.500 m, beam 4.350 m,
and draft 2.250 m. Natural periods for heave, pitch, and roll are 2.939, 3.429,
and 4.287 seconds, with damping ratios 0.45, 0.35, and 0.22.

The hull dimensions scale from the 28-foot preset by 42/28; all response
periods and horizontal time constants scale by the square root of that ratio.
Damping ratios are held fixed. This assumes geometric similarity, not a
particular production hull or loading condition.

Every plot includes all four incident-height cases: 0.27, 1.5, 4.0, and 8.5 m.
The first 60 seconds of each 20-minute, nominal 200 Hz record are plotted at
20 Hz. World acceleration excludes gravity; accelerometer specific force
includes gravity. Rotations and gyro readings describe the vessel rigid body.
These are analytical surrogate RAOs, not measured hull data. The largest sea
states extrapolate beyond the small-motion regime and are stress tests.
The CSV filenames continue to describe the incident wave parameters.

The complete dataset is \texttt{sim-data-files-vessel-rao-42ft.zip}; its internal
filenames and columns match the surface dataset. See \texttt{doc/vessel-rao.md}
for the response equations, coordinate conventions, and regular-wave details.
\clearpage
\section{Gerstner incident waves}
\VesselChart{gerstner_worldframe}{World-frame CG displacement, velocity and acceleration.}
\VesselChart{gerstner_imu_acc}{Body-frame accelerometer specific force.}
\VesselChart{gerstner_imu_gyro}{Body-frame angular velocity.}
\VesselChart{gerstner_euler}{Vessel Euler angles.}
\section{Jonswap incident waves}
\VesselChart{jonswap_worldframe}{World-frame CG displacement, velocity and acceleration.}
\VesselChart{jonswap_imu_acc}{Body-frame accelerometer specific force.}
\VesselChart{jonswap_imu_gyro}{Body-frame angular velocity.}
\VesselChart{jonswap_euler}{Vessel Euler angles.}
\section{Fenton incident waves}
\VesselChart{fenton_worldframe}{World-frame CG displacement, velocity and acceleration.}
\VesselChart{fenton_imu_acc}{Body-frame accelerometer specific force.}
\VesselChart{fenton_imu_gyro}{Body-frame angular velocity.}
\VesselChart{fenton_euler}{Vessel Euler angles.}
\section{Pmstokes incident waves}
\VesselChart{pmstokes_worldframe}{World-frame CG displacement, velocity and acceleration.}
\VesselChart{pmstokes_imu_acc}{Body-frame accelerometer specific force.}
\VesselChart{pmstokes_imu_gyro}{Body-frame angular velocity.}
\VesselChart{pmstokes_euler}{Vessel Euler angles.}
\section{Cnoidal incident waves}
\VesselChart{cnoidal_worldframe}{World-frame CG displacement, velocity and acceleration.}
\VesselChart{cnoidal_imu_acc}{Body-frame accelerometer specific force.}
\VesselChart{cnoidal_imu_gyro}{Body-frame angular velocity.}
\VesselChart{cnoidal_euler}{Vessel Euler angles.}
\end{document}
